\chapter{Technology Review}

\section{Overview}

This chapter provides a detailed examination of the technologies, tools,
frameworks, datasets, and algorithmic approaches that underpin the system
described in this dissertation. It is organised to mirror the architecture
of the project itself, moving from the AI generation layer through the
machine learning components, down to the backend infrastructure, database,
and frontend. Where possible, references are drawn from peer reviewed
publications, official standards bodies, and authoritative technical
documentation. The purpose of this chapter is to demonstrate that each
technology choice was informed by research and that the overall system
design is grounded in established principles from both software
engineering and applied machine learning. Each section tries to explain
what the technology is, why it was chosen for this project, and where
possible how it compares to the alternatives that were considered.

\section{Large Language Models and Generative AI}

\subsection{The Emergence of Large Language Models}

The capacity of this platform to generate quiz questions and coding
challenges on demand rests on large language models: neural networks
trained on enormous text corpora, capable of producing coherent output
in response to a natural language prompt. The architectural basis is
the Transformer introduced by Vaswani et al. \cite{vaswani2017attention},
which replaced the recurrent and convolutional structures of earlier
sequence models with a purely attention based mechanism. Attention
allows every position in a sequence to attend directly to every other
position, removing the sequential bottleneck that had limited previous
architectures and making it practical to train far larger models.

Subsequent work established that scaling these models produced
qualitative capability improvements. The bidirectional BERT model of
Devlin et al. \cite{devlin2019bert} showed that pre training on
unlabelled text followed by task specific fine tuning could set new
state of the art results across natural language understanding
benchmarks. The GPT line scaled autoregressive decoder only models
instead, and the work of Brown et al. \cite{brown2020gpt3} on GPT 3
demonstrated that scaling to 175 billion parameters produced few shot
and zero shot task performance, meaning these models could complete
tasks from a short description alone. This established the prompt
based interaction pattern this project relies on: a structured prompt
is sent to the model describing the desired output, and the model
generates a conforming response.

The project's primary backend is \texttt{gpt-4o-mini}, accessed through
the OpenAI API, a smaller, cost efficient member of the GPT 4 family
suited to reliable instruction following. A secondary backend using
Ollama \cite{ollama2025} with the \texttt{llama3.1:8b} model, descended
from Meta's LLaMA work \cite{touvron2023llama}, is also supported,
enabling locally hosted generation without reliance on an external API.
Abstracting both backends behind a common interface in
\texttt{ai\_models.py} means the system is not locked into any single
provider.

\subsection{Prompt Engineering}

Since LLMs are prompted rather than fine tuned in this project, output
quality depends heavily on prompt construction, a practice known as
prompt engineering \cite{liu2023promptsurvey}. The system uses two
prompt templates (\texttt{prompt\_guide.txt} for MCQ generation and
\texttt{coding\_prompt\_guide.txt} for coding challenges), each
specifying a strict JSON output schema to constrain the model to
produce parseable structured responses. This draws on the established
practice of using examples and formatting constraints in prompts, as
described by Brown et al. \cite{brown2020gpt3}. Recent work on chain
of thought prompting by Wei et al. \cite{wei2022chainofthought} has
shown that intermediate reasoning steps can improve output on multi
step tasks; the structured schema used here serves as a lighter weight
output scaffold.

A known limitation of prompt based integration is that LLMs do not
always produce strictly valid JSON. OpenAI has since introduced a
structured output mode \cite{openai2024structured} that enforces schema
conformance at the API level. This was not adopted at the time of
development but is identified as a clear improvement in the Conclusion.

\section{Machine Learning Foundations}

\subsection{Text Representation with TF IDF}

Both machine learning classifiers in this project use term frequency
inverse document frequency, or TF IDF, as the primary method of
transforming text into numerical feature vectors. TF IDF was introduced
by Sp\"{a}rck Jones \cite{sparck1972} in 1972 as a term weighting scheme
for information retrieval. The core insight is that a term's importance
to a document is proportional to how often it appears in that document,
but inversely proportional to how often it appears across the entire
corpus. Common words like \texttt{the} or \texttt{and} receive low
weights because they appear in almost every document and therefore carry
little discriminative information. Rare, domain specific terms receive
high weights because their presence is more informative. Salton and
Buckley \cite{salton1988tfidf} provided a systematic comparison of
alternative weighting schemes and established the form of TF IDF that is
now standard, and Ramos \cite{ramos2003tfidf} offers a more accessible
treatment suited to applied settings.

Formally, for a term $t$ in a document $d$ within a corpus $D$, the
TF IDF weight is defined as:
\[
\text{tf idf}(t, d, D) = \text{tf}(t, d) \times \log \left( \frac{|D|}{|\{d \in D : t \in d\}|} \right)
\]

This project uses scikit learn's \texttt{TfidfVectorizer}
\cite{pedregosa2011} with sublinear term frequency scaling, which
replaces the raw term frequency $\text{tf}(t, d)$ with $1 + \log(\text{tf}(t, d))$
to reduce the dominance of very frequent terms within a single document.
Two separate vectorisers are combined using \texttt{FeatureUnion}: a word
level vectoriser capturing unigrams through trigrams, and a character
level vectoriser operating on three to five character n grams. The
character level component is particularly valuable for code adjacent
text, where tokens such as \texttt{O(n)}, \texttt{dp[i]}, and
\texttt{arr[j]} form meaningful patterns that a word level tokeniser
would either split or discard. A word tokeniser typically treats
punctuation as a boundary and therefore cannot learn that \texttt{dp[i]}
is a recognisable motif, whereas a character n gram vectoriser will
naturally capture \texttt{dp[} and \texttt{]} as features.

\subsection{Logistic Regression}

The topic tag classifier uses logistic regression as its core estimator,
implemented within scikit learn \cite{pedregosa2011,buitinck2013api}.
Logistic regression is a discriminative linear classifier that models
the probability of class membership as a sigmoid function of a linear
combination of input features. For a feature vector $\mathbf{x}$ and
weight vector $\mathbf{w}$:
\[
P(y = 1 \mid \mathbf{x}) = \sigma(\mathbf{w}^T \mathbf{x} + b) = \frac{1}{1 + e^{-(\mathbf{w}^T \mathbf{x} + b)}}
\]

Parameters are estimated by maximising the log likelihood of the training
labels under this model, with L2 regularisation controlled by the
hyperparameter $C$ to prevent overfitting. The \texttt{class\_weight=}
\texttt{'balanced'} setting was used to adjust class weights inversely
proportional to their frequency in the training data, a standard
technique for addressing class imbalance \cite{japkowicz2002imbalance}.
Logistic regression was chosen over more complex models for this
component because TF IDF features tend to be high dimensional and
sparse, and linear classifiers are well known to perform competitively
in this regime \cite{hastie2009esl}. Deeper models such as transformer
based encoders would likely yield better absolute performance, but at
significantly higher inference cost, which would defeat the design goal
of sub second enrichment per question.

\subsection{Multi Label Classification}

Because a single coding problem can belong to multiple algorithmic
categories simultaneously (for example dynamic programming and binary
search together), topic tagging is inherently a multi label
classification problem \cite{zhang2014multilabel}. This project uses
scikit learn's \texttt{OneVsRestClassifier}, which trains one
independent binary classifier per label. The \texttt{MultiLabelBinarizer}
transforms the list of string tags into a binary vector of length
equal to the unique label count \cite{pedregosa2011}. An alternative,
the classifier chain of Read et al. \cite{read2011classifierchains},
was also evaluated but did not improve on the simpler one versus rest
approach on the LeetCode corpus.

The decision threshold was tuned per label on the validation set
rather than using a default 0.5 for every label. Lipton et al.
\cite{lipton2014thresholding} show that the F1 maximising threshold
depends on the positive class base rate and cannot be assumed constant.
The system uses a shrinkage regularised variant in which the per label
optimum is blended with a single global optimum to prevent overfitting
on rare labels where the validation sample is small.

\subsection{Gradient Boosting for Difficulty Prediction}

The difficulty classifier uses ensemble methods, comparing Random
Forest \cite{breiman2001randomforests}, HistGradientBoosting, and
Gradient Boosting estimators from scikit learn \cite{pedregosa2011}
with XGBoost \cite{chen2016xgboost}. Gradient boosting, as formalised
by Friedman \cite{friedman2001gradient}, builds an ensemble of shallow
decision trees by fitting each successive tree to the residuals of the
previous ensemble. The \texttt{HistGradientBoostingClassifier} bins
continuous features into discrete histograms before splitting,
dramatically reducing training time and handling missing values
natively. Linear SVMs \cite{cortes1995svm} were also evaluated as a
baseline. Model selection used five fold cross validation with
\texttt{RandomizedSearchCV} \cite{bergstra2012random}, scoring on F1
macro to ensure balance across difficulty classes. Random search is
preferred to grid search for larger hyperparameter spaces because, as
Bergstra and Bengio show, it explores more distinct values of the
important hyperparameters for the same computational budget.

\subsection{Proxy Labels and Class Imbalance}

The CodeNet dataset \cite{puri2021codenet} contains no human assigned
difficulty labels, so this project uses acceptance rate as a proxy,
binning problems into terciles: the top third labelled easy, the
middle medium, and the bottom hard. Using proxy labels is an
established technique in weakly supervised learning
\cite{zhou2017weaksupervision}. A noisy, cheap to obtain label is
substituted for an expensive authoritative one, and the model is
trained to recover the underlying signal despite the noise.

Class imbalance is a more serious issue for the topic classifier. The
LeetCode corpus is severely imbalanced: the Array tag appears on more
than a thousand problems while Depth First Search appears on fewer
than thirty. Japkowicz and Stephen \cite{japkowicz2002imbalance}
review the common mitigations, including cost sensitive learning,
oversampling, and undersampling. The \texttt{class\_weight='balanced'}
setting used in the logistic regression base classifier is an instance
of cost sensitive learning. Synthetic oversampling via SMOTE
\cite{chawla2002smote} was considered but not adopted, as its
interaction with high dimensional sparse TF IDF vectors is not
straightforwardly meaningful.

\section{Datasets}

\subsection{LeetCode Dataset for Topic Tag Classification}

The topic tag classifier was trained on a LeetCode derived dataset obtained from Kaggle \cite{kaggleleetcode2025}, distributed as a CSV file containing problem titles, descriptions in HTML format, and associated topic tags. LeetCode is a widely used competitive programming and interview preparation platform, and its problems have well established community assigned tags such as Array, Hash Table, Dynamic Programming, and Two Pointers. After cleaning, parsing multi value tag strings, and filtering out tags appearing in fewer than sixty examples (to remove labels with insufficient training signal), this dataset provided a realistic and domain appropriate training corpus for the tagging task. An initial exploratory analysis on the filtering threshold established that too low a threshold admitted labels where the classifier could not learn any useful pattern, while too high a threshold collapsed the problem into a handful of very common tags that were already trivially easy to predict.

\subsection{IBM Project CodeNet}

The difficulty classifier was trained on IBM Project CodeNet
\cite{puri2021codenet}, published by Puri et al. at NeurIPS 2021.
CodeNet is a large scale dataset containing approximately fourteen
million code samples across 4,053 competitive programming problems in
fifty five programming languages. Each problem submission is annotated
with metadata including its acceptance status, CPU runtime, memory
usage, and code size. The dataset was constructed from two online
judging platforms and has been used widely in research on AI for code.

For this project, per problem statistics were aggregated from
individual submission CSV files, and problem descriptions were parsed
from HTML to extract lexical and structural features. The scale and
richness of CodeNet's metadata made it the most suitable publicly
available dataset for engineering the numeric and textual features
required by the difficulty model. The dataset is released with a
permissive licence, which was a practical consideration for a final
year project: no restrictive terms constrain how the derived model
artefacts are used or distributed.

\section{Backend Framework: FastAPI}

The backend is built using FastAPI \cite{fastapi2025}, a Python web
framework released in 2018 by Sebasti\'{a}n Ram\'{i}rez, layered on
Starlette \cite{starlette2025} for asynchronous request handling and
Pydantic \cite{pydantic2025} for data validation. Routes declare the
types of their parameters and return values using Python type hints,
which FastAPI uses to automatically validate incoming requests,
serialise responses, and generate interactive OpenAPI documentation.
Independent benchmarks place FastAPI among the fastest Python
frameworks, comparable in throughput to Node.js and Go when run under
the Uvicorn ASGI server \cite{fastapi2025}.

FastAPI was chosen over Flask and Django REST Framework for two
reasons. Its native async support allows the backend to handle
concurrent AI API calls without blocking, which matters because AI
generation latency dominates overall request time. Automatic Pydantic
validation removes a large category of defensive checking code from
route handlers. Lathkar's book
\cite{lathkar2023} provides a comprehensive treatment of the
framework's async architecture. All request and response structures
are defined as Pydantic \texttt{BaseModel} subclasses, so malformed
input is caught at the API boundary and never reaches business logic,
a pattern of edge validation that reduces defensive code and eases
diagnosis of failures \cite{martin2008cleancode}.

\section{Database: PostgreSQL}

PostgreSQL \cite{postgres2025} is used as the relational database,
accessed through \texttt{psycopg} version 3 with a \texttt{dict\_row}
factory. PostgreSQL is a mature, open source, object relational
database with roots in the POSTGRES project at UC Berkeley
\cite{stonebraker1986postgres}, ACID compliant since 2001 and
implementing multi version concurrency control. The schema stores
user accounts, profiles, quiz results, experience and levelling data,
and the social graph of friendships. A relational database was chosen
over a document store because the social features involve many to many
relationships with status qualified edges, which benefit from
normalised modelling and join based querying. Native PostgreSQL
features used include \texttt{TEXT[]} array columns (for topic tags on
quiz results), \texttt{JSONB} columns (for full quiz session content),
and \texttt{CHECK} constraints (to enforce at the database level that
a user cannot friend themselves).

TablePlus \cite{tableplus2025} was used throughout development as the primary
graphical interface for inspecting and managing the PostgreSQL database. As a
native GUI client, it supports direct connection to PostgreSQL and provides
inline editing of rows and table structure, as well as advanced filtering to
quickly locate specific records. During development, TablePlus served as the
main tool for verifying schema correctness, manually inspecting table contents after test runs, and executing ad hoc queries to debug application behaviour. Its built-in SQL editor with syntax highlighting and the ability to split query results across multiple tabs made it straightforward to cross-reference data across the user accounts, quiz results, and friendships tables simultaneously. The code review feature and snapshot capability also provided a safety net when making direct modifications to the development database, ensuring unintended changes could be identified before being committed.

\section{Containerisation: Docker and Docker Compose}

The entire application is containerised using Docker
\cite{merkel2014docker}. Docker packages applications and their
dependencies into lightweight containers that run as host processes
isolated via Linux kernel namespaces and control groups. Its key
promise, that a containerised application behaves identically wherever
it runs, matters for a project of this complexity, where the backend
depends on a specific Python version, Node.js for parts of the code
execution engine, and pre trained model artefacts. Docker Compose
\cite{dockercompose2025} orchestrates multiple containers
declaratively in a single YAML file. Three services are defined:
PostgreSQL, the Python backend, and the React frontend served by
Nginx. The backend Dockerfile installs Python 3.11 slim with Node.js
22 alongside for the JavaScript and TypeScript runtimes. The frontend
uses a two stage build: Node 22 alpine produces the Vite artefact,
which is served by Nginx stable alpine with single page fallback
routing, gzip compression, and one year cache headers on hashed
assets.

\section{Frontend: React and TypeScript}

The frontend is built using React \cite{react2013} and TypeScript
\cite{typescript2025}, compiled by Vite \cite{vite2025}. React is an
open source JavaScript library for component based interfaces,
developed at Meta and open sourced in 2013. Its central innovation is
the virtual DOM: rather than directly manipulating the browser DOM on
every state change, React maintains an in memory UI representation
and computes the minimal set of changes required to synchronise the
real DOM with the current state. TypeScript adds a static type system
to JavaScript, catching type errors at compile time. For a project of
this size, where data flows from an AI model through a parser, a
machine learning enrichment layer, and a REST API before reaching the
frontend, type safety at the boundary protects against category
mismatches: a Pydantic schema on the backend and a matching
TypeScript interface on the frontend form a contract the compiler can
check. Vite provides near instant hot module replacement and produces
optimised production bundles. State is managed using
React hooks and an \texttt{AuthContext} provider that exposes the JWT
and user session. The Monaco Editor \cite{monaco2025}, the same engine
that powers Visual Studio Code, is embedded in the coding sandbox to
provide syntax highlighting, auto indentation, and language aware
features for all five supported languages.

\section{Authentication: JSON Web Tokens and bcrypt}

User authentication uses JSON Web Tokens as standardised in RFC 7519
\cite{jwt2015}. A JWT is a compact URL safe representation of claims
encoded as JSON and signed with HMAC SHA256. On login, the backend
issues a signed token carrying the user id and an expiry timestamp;
the client stores it and includes it in the Authorization header on
subsequent requests, where the backend verifies the signature without
a database lookup. This stateless pattern scales well because no
server side session state is held. The PyJWT library handles
encoding and decoding on the backend. Passwords are hashed with
bcrypt \cite{provos1999bcrypt}, an adaptive hashing function based on
Blowfish whose configurable cost parameter keeps it secure as
attacker hardware improves, unlike a bare cryptographic hash such as
SHA 256 which is fast by design and therefore easy to brute force.
OWASP authentication guidance \cite{owasp2021} recommends adaptive
hashing functions such as bcrypt, scrypt, or argon2.

\section{Code Execution Engine}

The code execution engine supports Python, JavaScript, TypeScript,
Java, and C\#. For each submission, the backend writes the user's
code to a temporary file and invokes the runtime via
\texttt{subprocess.run} with a five second timeout. The backend Docker
image installs Node.js 22 to support JavaScript and TypeScript, and
both Java and C\# runtimes are also provided. For test case
evaluation, the engine compares actual standard output against the
expected output per case. Subprocess isolation is a common pattern in
online judges; it does not match kernel level isolation systems such
as Isolate, but is appropriate for the scope and threat model of this
project, where users are authenticated and have incentive to learn
rather than to subvert the platform. Production deployment with
untrusted users would need stronger isolation, flagged as future work.

\section{Data Interchange and Continuous Integration}

All frontend and backend communication uses JSON (RFC 7159
\cite{json2014}), consistent with the REST architectural style
\cite{fielding2000rest}. FastAPI automatically generates an OpenAPI
3.0 specification \cite{openapi2025} from the Pydantic model
definitions, documenting every endpoint's request and response
schemas. The continuous integration pipeline uses GitHub Actions
\cite{githubactions2025} with workflows defined in
\texttt{.github/workflows/}, running three parallel jobs covering the
frontend (ESLint, TypeScript type checking, Vite build), the backend
(Black, Flake8, Python syntax validation), and the machine learning
models directory.

\section{Scientific Computing Libraries}

Beyond scikit learn \cite{pedregosa2011,buitinck2013api}, the machine
learning components use NumPy \cite{numpy2020} for array operations,
pandas \cite{mckinney2010pandas} for tabular data manipulation, and
joblib \cite{joblib2025} for serialising trained pipelines to
\texttt{.pkl} files. The \texttt{ColumnTransformer} and
\texttt{FeatureUnion} classes build modular preprocessing pipelines
that apply different transformations to different feature subsets and
concatenate the resulting matrices before model fitting. This pipeline
based approach ensures preprocessing is consistently applied at
inference time using the parameters fitted on training data,
eliminating training and serving skew. Hastie, Tibshirani, and
Friedman \cite{hastie2009esl} provide the statistical learning
foundation for the estimators used, and Goodfellow, Bengio, and
Courville \cite{goodfellow2016deep} offer a reference should future
iterations replace TF IDF with a learned representation.